\section{Heilsgeschichte}
        
    \begin{block}[Erklärung]
        Die folgenden Ausführungen habe ich zum grössten Teil von Dr. Michael J. Vlach. Er ist Dozent in Kalifornien und hat sich jahrzehntelang mit diesem Thema beschäftigt.

        Die Heilsgeschicht in der Bibel ist ein zentrales Thema. Sie ist der rote Faden durch das Wort Gottes. Nur wenn wir die Heilsgeschichte verstehen, können wir auch den Plan Gottes mit uns verstehen.
  \end{block}
  
   \begin{block}[Das Reich Gottes als Thema in der Schrift]
    Wenn man die Bibel von Anfang bis Ende durchliest, was fällt einem da auf?
   \begin{itemize}
    \item Die Bibel erzählt über das Heil des Menschen
    \item über die Errettung des Menschen
    \item Geschichte eines Volkes
    \item Geschichte von Jesus
    \item Die Gründung der Gemeinde
    \item Prophetische Texte
   \end{itemize}
   Alles richtig. Aber wieso hat Gott den Menschen überhaupt erschaffen? Der Mensch, den er geschaffen hat, war ja eigentlich sehr gut und er hat den Menschen in das Paradies gestellt? Wieso eine Heilsgeschichte? Viele der vorher genannten Themen sind die Folgen des Sündenfalls und diesen Sündenfall war von Gott so nicht gewollt. 
   \end{block}
   Ein grosses zentrales Thema ist das Reich Gottes. Bereits im Schöpfungsbericht zeigt sich der Herr als der regierende König, über seine zuvor geschaffene Welt \biblezit{IMos}{1:2:} Auch der Mensch sollte in diesem vom Gott erschaffnen Reich, eine Schlüsselstellung einnehmen. Zu lesen in den Versen
    \begin{bibelbox}{SCHL}{1Mos}{1:26-28}
      \hh{26}Und Gott sprach: Lasst uns Menschen machen nach unserem Bild, uns ähnlich; die sollen herrschen über die Fische im Meer und über die Vögel des Himmels und über das ganze Vieh und über die ganze Erde, auch über alles Gewürm, das auf der Erde kriecht!\\
      \hh{27}Und Gott schuf den Menschen in seinem Bild, im Bild Gottes schuf er ihn; als Mann und Frau schuf er sie.\\
      \hh{28}Und Gott segnete sie; und Gott sprach zu ihnen: Seid fruchtbar und mehrt euch und füllt die Erde und macht sie euch untertan; und herrscht über die Fische im Meer und über die Vögel des Himmels und über alles Lebendige, das sich regt auf der Erde!
    \end{bibelbox}
   Dann aber wurde der Mensch zum Sünder wie in Kapitel 3 in den Versen 6 und 7 beschrieben ist. 
   \begin{bibelbox}{SCHL}{IMos}{3:6-7}
    \hh{6}Und die Frau sah, dass von dem Baum gut zu essen wäre, weil er weise macht; und sie nahm von seiner Frucht und ass, und sie gab davon auch ihrem Mann, der bei ihr war.\\
    \hh{7}Da wurden ihnen beiden die Augen geöffnet, und sie erkannten, dass sie nacht waren; und sie banden sich Feigenblätter um und machten sich Schurze.
   \end{bibelbox}
   Diese Sünde überwindet Gott, indem er bereits vor Grundlegung der Welt geplant hat, wie die Sünde überwunden und das Reich
   Gottes trotzdem aufgerichtet werden kann. Dieser Stufenplan, welcher sich über mehrere Jahrtausende hinzieht, wird in der Bibel detailliert beschrieben. Das Ziel Gottes besteht darin, alle ihm ergebenen Menschen wie ursprünglich vorgesehen als Verwalter über seine Schöpfung zu stellen. So lesen wir am Ende der Bibel
   \begin{bibelbox}{SCHL}{Offb}{22:3-5}
    \hh{3}Und es wird keinen Fluch mehr geben; und der Thron Gottes und des Lammes wird in ihr sein, und seine Knechte werden ihm dienen;\\
    \hh{4}und sie werden sein Angesicht sehen, und sein Name wird auf ihren Stirnen sein\\
    \hh{5}Und es wird dort keine Nacht mehr geben, und sie bedürfen nicht eines Leuchters, noch des Lichtes der Sonne, denn Gott, der Herr, erleuchtet sie; und sie werden herrschen von Ewigkeit zu Ewigkeit.
   \end{bibelbox}
   Von der ersten bis zur letzten Seite der Bibel, lesen wir wie Gott Schritt für SChritt zum Ziel kommt, mit uns Menschen das Reich Gottes aufzurichten.
   Vlach teilt die Entfaltung des Reiches Gottes in 5 Phasen ein und wir werden die einzelnen Phasen in den nächsten Gottesdiensten anschauen:
   
   \begin{enumerate}
    \item Schöpfung \bibleverse{IMos}(1-2:)\\
     Der Mensch wird als Verwalter einer vollendeten Schöpfung geschaffen. Gott der König 
beauftragt die Menschen, in seinem Namen über die Schöpfung zu herrschen (1Mo 
1,26-28). 
    \item Sündenfall \bibleverse{IMos}(3:)\\
    Der Sündenfall zeigt über das menschliche Versagen, dass der Mensch nicht willens ist, 
die Aufgabe als Herrscher über die Erde anzunehmen. Die Schöpfung kommt dadurch 
unter den Einfluss Satans. Beides, die Schöpfung und die Menschen fallen unter den 
Fluch der Sünde
    \item Verheissung 1. Mose 3,15 - Maleachi\\
    Gott verheißt, dass der „Same“ der Frau letztendlich über das Böse siegen wird (1Mo 
3,15). Damit wird der Sündenfall überwunden werden, und auch der Mensch wird 
seiner ursprünglichen Berufung entsprechend über die Schöpfung herrschen können.
    \item Ankunft des Königs Evangelien Briefe\\
     Über seinen Tod am Kreuz erfüllt Jesus Christus alle Bedingungen, die zur Aufrichtung 
des Reiches Gottes nötig sind. Zugleich werden die Gläubigen dieser geistlichen 
Realität in direkter Weise teilhaftig. 
    \item Wiederherstellung Offenbarung\\
     Mit der heute noch ausstehenden Wiederkunft Jesu, wird das Reich Gottes auch 
physisch vollends aufgerichtet werden. Dieses Reich wird am Ende seinerseits durch ein 
 vollkommenes Reich abgelöst, in welchem die Erlösten die ursprüngliche Absicht 
Gottes freiwillig und mit großer Freude wahrnehmen werden.
   \end{enumerate}

Wenn wir beim Bibellesen nur einzelne Verse aussuchen, werden wir nie diesen roten Faden finden. Wenn wir aber die Bibel von Anfang bis Ende durchlesen, können wir schön sehen, wie die einzelnen Bücher ineinander eingehen. Und das ist einer der Wunder dieses Buches.

Menschen können nicht über Jahrhunderte sogar Jahrtausende ein Buch schreiben, dass vollständig in sich stimmig ist. Verstehen können wir das nur, wenn wir Jesus in und mit uns haben. Darum singen wir nun gemeinsam das Lied. Nur Christus in mir.
